\documentclass[10pt,a4paper]{article}
\usepackage[T1]{fontenc}
\usepackage[utf8]{inputenc}
\usepackage[french]{babel}
\usepackage[a4paper,margin=1.25cm]{geometry}
\usepackage[hidelinks]{hyperref}
\usepackage{enumitem}
\usepackage{titlesec}
\usepackage{xcolor}
\definecolor{graytext}{gray}{0.35}
\definecolor{accent}{HTML}{1A4F8B}
\pagestyle{empty}
\setlength{\parindent}{0pt}
\setlist[itemize]{leftmargin=1em,itemsep=0.5pt,topsep=1pt,parsep=0pt}
\titleformat{\section}{\normalsize\bfseries\color{accent}}{}{0pt}{}[\vspace{-3pt}\rule{\linewidth}{0.4pt}]
\titlespacing*{\section}{0pt}{6pt}{6pt}
\newcommand{\cventry}[3]{\textbf{#1}\ \textcolor{graytext}{· #2}\hfill\textcolor{graytext}{\small #3}\par}
\newcommand{\tagline}[1]{\textcolor{graytext}{\small #1}\par}
\begin{document}

{\LARGE\bfseries TOUNEKTI Mohamed Rafaa}\hfill{\small DevOps \& DevSecOps}\par
\vspace{3pt}
{\footnotesize \href{https://linkedin.com/in/mohamed-rafaa-tounekti}{linkedin.com/in/mohamed-rafaa-tounekti} \textcolor{graytext}{·} \href{tel:+5511929017154}{+55 11 92901-7154} \textcolor{graytext}{·} \href{mailto:rafaabou@gmail.com}{rafaabou@gmail.com} \textcolor{graytext}{·} Basé à São Paulo (CPF Ready) – Brésil \textcolor{graytext}{·} \href{https://github.com/rafaabou}{github.com/rafaabou} \textcolor{graytext}{·} \href{https://devops-devsecops.kdns.fr/}{devops-devsecops.kdns.fr}}\par
\vspace{2pt}\hrule\vspace{2pt}
\section*{Profil}
Ingénieur DevOps et DevSecOps spécialisé dans la construction de pipelines CI/CD fiables, la sécurisation des clusters Kubernetes et l'automatisation de l'infrastructure à grande échelle. Je travaille avec AWS, les systèmes Linux et les environnements conteneurisés, avec une expérience approfondie de GitLab CI/CD, Terraform, Ansible et des stacks de monitoring.\par
\section*{Compétences Techniques}
\textbf{Cloud}: AWS (EC2, EKS, ECR, IAM, S3, VPC, Lambda)\par
\textbf{Conteneurs \& Orchestration}: Docker, Kubernetes, Helm, Kustomize, Istio, OPA, Kyverno\par
\textbf{CI/CD \& GitOps}: GitLab CI/CD, Jenkins, ArgoCD, Flux CD, GitLab Runners\par
\textbf{Infrastructure en tant que Code}: Terraform, Ansible\par
\textbf{Sécurité}: Trivy, GitLeaks, Semgrep, OWASP ZAP, SonarQube, Checkov, KubeLinter, Falco, HashiCorp Vault, AWS Secrets Manager\par
\textbf{Monitoring \& Logging}: Prometheus, Grafana, Loki, Alertmanager\par
\textbf{Réseau · Scripts · Bases de données · Stockage}: Nginx Ingress, Traefik, Cert-Manager, External-DNS, Bash, Python, PostgreSQL, MongoDB, Redis, NFS, Harbor, Nexus\par
\textbf{Virtualisation · Outils IA}: KVM, VirtualBox, Vagrant, MCP, OpenCode, AI-assisted infra\par
\section*{Expérience}
\cventry{Ingénieur DevOps \& DevSecOps}{TUNISIE SOFTWARE FACTORY \& BE-SOFTILYS}{Juin 2025 – Octobre 2024}
\begin{itemize}
  \item Construit des pipelines GitLab CI/CD avec scan Trivy, GitLeaks et Semgrep. Images vulnérables bloquées avant le staging
  \item Déployé une sécurité Zero-Trust sur les clusters Kubernetes avec Istio mTLS, NetworkPolicies et OPA Gatekeeper
  \item Intégré HashiCorp Vault dans les workflows de déploiement avec rotation automatique des secrets via External Secrets Operator
  \item Mis en place des tableaux de bord Prometheus \& Grafana avec routage Alertmanager ; le MTTR a nettement diminué
  \item Configuré HPA et PodDisruptionBudgets pour gérer les pics de trafic sans perdre de requêtes
  \item Automatisé le provisionnement des certificats TLS avec Cert-Manager pour toutes les ressources Ingress
  \item Effectué des tests de validation d'infrastructure de bout en bout après chaque cycle de provisionnement
  \item Appliqué des contrôles de sécurité CI/CD avec des seuils de vulnérabilité ; pipelines bloqués au-delà des seuils
  \item Centralisé le suivi des vulnérabilités dans DefectDojo pour le suivi des corrections
  \item Débogué des problèmes d'exécution Kubernetes — CrashLoopBackOff, ImagePullBackOff, échecs d'ordonnancement
  \item Rédigé des runbooks et des documents de déploiement ; rapports de validation conservés
  \item Discussions pré-vente avec des clients potentiels sur les options d'architecture
  \item Configuré des serveurs MCP pour les assistants IA afin d'interroger les données Kubernetes, CI/CD et monitoring
  \item Intégré des assistants IA dans les workflows quotidiens pour l'analyse de logs, le triage d'incidents, le diagnostic
  \item Utilisé des outils IA pour accélérer le développement Terraform et Ansible ; moins d'erreurs de configuration
  \item Règles de protection de branche avec revues obligatoires et vérifications automatisées des pipelines
  \item Connecté Alertmanager à Jira ; les alertes génèrent automatiquement des tickets
  \item Sécurité de la chaîne d'approvisionnement : génération SBOM, signature Cosign, suivi des dépendances
  \item Déployé Falco pour la sécurité runtime sur les charges de travail Kubernetes
\end{itemize}
\tagline{Kubernetes · Docker · Helm · ArgoCD · GitLab CI/CD · Terraform · AWS · GitLeaks · Semgrep · Trivy · OPA Gatekeeper · Kyverno · DefectDojo · Prometheus · Grafana · Istio · HashiCorp Vault · Harbor}
\vspace{1pt}
\vspace{3pt}
\cventry{Ingénieur DevOps}{DOPIFY SOFTWARE \& TECHNOLOGIES}{Avril 2024 – Janvier 2024}
\begin{itemize}
  \item Conçu des pipelines CI/CD pour plusieurs équipes avec SAST, SCA et scan de secrets dans chaque build
  \item Automatisé le provisionnement des serveurs en dev, staging, production avec Ansible ; éliminé la dérive manuelle
  \item Mis en place Flux CD pour une infrastructure GitOps ; environnements déclaratifs et auditable
  \item Déployé Grafana Loki pour l'agrégation centralisée des logs ; analyse des causes plus rapide
  \item Déploiements sans temps d'arrêt avec des stratégies blue-green et canary via Kustomize
  \item Configuré Nginx Ingress Controller avec terminaison TLS pour le routage externe
  \item Provisionné des runners GitLab auto-hébergés avec exécuteur Docker ; les temps d'attente CI ont diminué
  \item Alertmanager a routé les incidents vers Slack et email ; les bonnes personnes contactées au bon moment
  \item Automatisé des scripts de test pour valider l'état de l'infrastructure après provisionnement
  \item Géré la gouvernance des pipelines et la planification des transitions pour les releases de production
  \item Identifié les dégradations de performance via les logs système et application
  \item Documenté l'architecture d'infrastructure, les procédures de déploiement, les problèmes connus
  \item Protection de branche Git avec revues obligatoires et vérifications de statut ; contrôles qualité avant merge
\end{itemize}
\tagline{AWS · Kubernetes · Docker · Terraform · Ansible · Flux CD · Grafana · Prometheus · GitLab CI/CD · HashiCorp Vault · Trivy · SonarQube · Kustomize · Nginx · Alertmanager}
\vspace{1pt}
\vspace{3pt}
\cventry{Ingénieur DevOps}{GEO-TRACKING GROUP}{Octobre 2023 – Avril 2022}
\begin{itemize}
  \item Reconstruit les pipelines CI/CD ; temps de déploiement réduit de jours à heures pour 15+ microservices
  \item Réduit la taille des images Docker de 60\% avec des builds multi-étapes ; coûts de pull et de stockage réduits
  \item Provisionné l'infrastructure AWS avec des modules Terraform ; versionné, reproductible
  \item Intégré SonarQube pour les vérifications automatisées de qualité de code et le suivi de la dette technique
  \item Rédigé des règles d'alerte Prometheus pour la détection proactive ; les pannes non planifiées ont chuté de 40\%
  \item Automatisé les mises à jour de dépendances et les correctifs de sécurité via GitOps sur 200+ dépendances
  \item Géré les images Docker et les manifests Kubernetes dans Nexus comme source unique de vérité
  \item Automatisé les sauvegardes de bases de données avec restauration ponctuelle pour PostgreSQL et MySQL
  \item Renforcé la sécurité des pods avec SecurityContext et segmentation réseau via IPTables/UFW
  \item Construit des bibliothèques Jenkins partagées réutilisables pour standardiser les templates de pipelines
  \item Effectué des tests d'installation et de mise à niveau ; résultats documentés et écarts signalés
  \item Surveillé la capacité de stockage et les métriques de performance ; configurations ajustées pour la disponibilité
  \item Répondu aux incidents et rédigé des rapports post-incident avec mesures correctives
  \item Intégré LDAP pour l'authentification centralisée
\end{itemize}
\tagline{Jenkins · Docker · Kubernetes · Terraform · Ansible · ArgoCD · SonarQube · Nexus · Prometheus · Grafana · Python · Bash · AWS · PostgreSQL · Cert Manager}
\vspace{1pt}
\vspace{3pt}
\cventry{Administrateur Système}{DIVERSES ORGANISATIONS}{Janvier 2022 – Février 2016}
\begin{itemize}
  \item Géré des serveurs Linux (RHEL, Ubuntu, Debian) — correctifs, mises à jour sécurité, optimisation
  \item Automatisé le provisionnement de VM avec Ansible ; temps de configuration d'environnement réduit
  \item Administré des routeurs, commutateurs, pare-feux ; maintenu les politiques réseau
  \item Mis en place le partage de fichiers NFS pour l'accès centralisé au stockage entre équipes
  \item Configuré DNS, DHCP, services réseau internes
  \item Conçu et testé des procédures de reprise après sinistre avec validation régulière des sauvegardes
  \item Géré MySQL et PostgreSQL — sauvegarde, restauration, optimisation des performances
  \item Configuré les règles de pare-feu IPTables et UFW pour segmenter le réseau et restreindre l'accès
  \item Déployé des serveurs web Apache et Nginx pour des applications internes/externes
  \item Intégré Active Directory pour l'authentification centralisée sur Linux/Windows
  \item Géré les comptes utilisateurs, clés SSH, politiques sudo dans des environnements d'entreprise
  \item Fourni un support L2/L3 avec analyse des causes racines
  \item Aidé les utilisateurs finaux avec des problèmes matériels, logiciels, réseau ; sessions de formation
  \item Maintenu la documentation des processus IT, configurations réseau, dépannage
  \item Configuré des systèmes de stockage sur IPv4 avec adressage et routage appropriés
  \item Collaboré avec les équipes techniques pour planifier les améliorations et l'automatisation
  \item Automatisé les workflows alerte-ticket avec résolution gérée par SLA
\end{itemize}
\tagline{Linux · DHCP · DNS · Active Directory · MySQL · PostgreSQL · VMware · KVM · Ansible · Bash · VPN · Firewalls · NFS · Apache · Nginx}
\vspace{1pt}
\section*{Formation}
\textbf{Licence}\par
\tagline{Technologies de l'Information (Réseaux et Services Informatiques) · I.S.E.T. · 2011 – 2014}
\textbf{Baccalauréat}\par
\tagline{Science Informatique · 2010 – 2011}
\section*{Langues}
Arabe (Natif), Français (Professionnel), Anglais (Professionnel)\par
\section*{Réalisations Clés}
\begin{itemize}
  \item Architecture de sécurité Zero-Trust sur Kubernetes
  \item Scan de sécurité automatisé dans les pipelines CI/CD
  \item GitOps avec ArgoCD \& Flux CD
  \item Haute disponibilité avec PDB et auto-scaling
  \item Stack complet de monitoring et d'alerting
\end{itemize}
\vspace{6pt}\hrule\vspace{3pt}
{\footnotesize\color{graytext} Analyse de logs assistée par IA · triage d'incidents · diagnostic d'infrastructure · serveurs MCP\par}
\end{document}